\uuid{g5Ka}
\titre{ Forward pass et interprétation d'un neurone sigmoïde }
\niveau{L3}
\module{Analyse de données}
\chapitre{Réseaux de neurones}
\sousChapitre{Forward pass, fonction sigmoïde}
\theme{Neurone sigmoïde, entropie croisée binaire, degré de confiance}
\auteur{Maxime Nguyen}
\datecreate{2026-03-19}
\organisation{AMSCC}
\difficulte{2}

\contenu{
\texte{
  On considère un réseau réduit à un seul neurone sigmoïde, de paramètres $w = 2$ et $b = -1$,
  recevant une entrée $x = 1$. On rappelle que $\sigma(z) = \frac{1}{1+e^{-z}}$ et que $\sigma(1) \approx 0{,}73$.
}
\begin{enumerate}

\item
  \question{Calculer $z = wx + b$, puis $\hat{p} = \sigma(z)$.}
  \reponse{
    On a $z = 2 \times 1 + (-1) = 1$, donc $\hat{p} = \sigma(1) \approx 0{,}73$.
  }

\item
  \question{La vraie classe est $y = 0$. Le modèle se trompe-t-il ? Avec quel degré de confiance prédit-il la classe $1$ ?}
  \reponse{
    Le modèle prédit $\hat{p} \approx 0{,}73$, soit une probabilité de $73\,\%$ pour la classe $1$.
    Or la vraie classe est $y = 0$ : le modèle se trompe, avec un degré de confiance élevé ($73\,\%$) dans la mauvaise direction.
  }

\item
  \question{Sans calculer la perte d'entropie croisée binaire (BCE) exacte, anticiper si elle sera plutôt faible ou forte. Justifier.}
  \reponse{
    La BCE vaut $-\log(1 - \hat{p}) \approx -\log(0{,}27) \approx 1{,}31$.
    Elle est \emph{forte} : le modèle est confiant ($\hat{p} \approx 0{,}73$) mais dans la mauvaise classe.
    Plus le modèle est confiant et erroné, plus $-\log$ d'une probabilité proche de $0$ est grande.
  }

\end{enumerate}
}
